\documentclass[12pt,a4paper]{article}

% ==================== 基础包 ====================
\usepackage[UTF8]{ctex}
\usepackage[margin=2.5cm]{geometry}
\usepackage{graphicx}
\usepackage{booktabs}
\usepackage{amsmath,amssymb,amsfonts}
\usepackage{listings}
\usepackage{xcolor}
\usepackage{hyperref}
\usepackage{float}
\usepackage{caption}
\usepackage{subcaption}
\usepackage{enumitem}
\usepackage{algorithm}
\usepackage{algpseudocode}
\usepackage{biblatex}
\addbibresource{references.bib}

% ==================== 代码样式 ====================
\lstset{
    language=Python,
    basicstyle=\ttfamily\small,
    keywordstyle=\color{blue}\bfseries,
    commentstyle=\color{green!60!black},
    stringstyle=\color{red!80!black},
    numbers=left,
    numberstyle=\tiny\color{gray},
    stepnumber=1,
    numbersep=8pt,
    backgroundcolor=\color{gray!10},
    frame=single,
    rulecolor=\color{gray!30},
    breaklines=true,
    breakatwhitespace=true,
    showstringspaces=false,
    tabsize=4,
    captionpos=b,
    xleftmargin=1.5em,
    framexleftmargin=1em
}

% ==================== 超链接设置 ====================
\hypersetup{
    colorlinks=true,
    linkcolor=blue,
    filecolor=magenta,
    urlcolor=cyan,
    citecolor=green!50!black
}

% ==================== 标题信息 ====================
\title{\textbf{机器学习导论大作业报告} \\ \vspace{0.5em} \large CIFAR-10图像分类实验}
\author{
    \begin{tabular}{rl}
        \textbf{专业:} & 计算机科学与技术 \\
        \textbf{班级:} & 计科2101 \\
        \textbf{学号:} & 20210001 \\
        \textbf{姓名:} & 张三 \\
    \end{tabular}
}
\date{\today}

\begin{document}

\maketitle

\begin{abstract}
本报告系统研究了CIFAR-10数据集上的图像分类问题，分别实现了多层感知机(MLP)、卷积神经网络(CNN)和ResNet-18三种模型架构。通过大量消融实验，深入分析了Batch Normalization、Dropout、优化器选择、学习率调度和数据增强等关键技术对模型性能的影响。实验结果表明，ResNet-18在测试集上达到了95.16\%的准确率，显著优于CNN(83.66\%)和MLP(39.31\%)。本报告详细记录了实验设计、结果分析和结论，为深度学习图像分类任务提供了有价值的参考。
\end{abstract}

\tableofcontents
\newpage

% ==================== 第1章 ====================
\section{问题定义与理解}

\subsection{任务描述}

CIFAR-10是由Alex Krizhevsky等人收集整理的一个广泛用于图像分类任务的标准数据集\cite{krizhevsky2009learning}。该数据集包含10个不同类别的RGB彩色图像，每个图像的尺寸为$32 \times 32$像素。这10个类别分别是：飞机(airplane)、汽车(automobile)、鸟类(bird)、猫(cat)、鹿(deer)、狗(dog)、青蛙(frog)、马(horse)、船(ship)和卡车(truck)。

\subsection{数据集规模}

CIFAR-10数据集的具体规模如下：
\begin{itemize}[leftmargin=2em]
    \item \textbf{训练集:} 50,000张图像，每个类别5,000张
    \item \textbf{测试集:} 10,000张图像，每个类别1,000张
    \item \textbf{图像尺寸:} $32 \times 32 \times 3$ (RGB三通道)
    \item \textbf{类别数:} 10个互斥类别
\end{itemize}

\subsection{评估指标}

本实验采用\textbf{分类准确率}(Classification Accuracy)作为主要评估指标，定义为：
\begin{equation}
    \text{Accuracy} = \frac{\text{正确分类的样本数}}{\text{总样本数}} \times 100\%
\end{equation}

此外，我们还记录了每个类别的准确率、混淆矩阵以及训练过程中的损失曲线，以全面评估模型性能。

\subsection{任务难点分析}

CIFAR-10图像分类任务面临以下挑战：
\begin{enumerate}[leftmargin=2em]
    \item \textbf{低分辨率:} $32 \times 32$的图像尺寸导致细节信息有限
    \item \textbf{类间相似性:} 某些类别(如猫和狗、鹿和马)在视觉上非常相似
    \item \textbf{类内差异性:} 同一类别的物体可能呈现不同姿态、颜色和背景
    \item \textbf{小样本量:} 每类仅5,000张训练图像，容易导致过拟合
\end{enumerate}

% ==================== 第2章 ====================
\section{数据分析及处理}

\subsection{数据探索}

在构建模型之前，我们首先对CIFAR-10数据集进行了详细的数据探索分析。

\subsubsection{类别分布}

CIFAR-10数据集的类别分布是完全均衡的，如图~\ref{fig:class_distribution}所示。训练集中每个类别恰好包含5,000张图像，测试集中每个类别恰好包含1,000张图像。这种均衡分布避免了类别不平衡带来的偏差问题。

\begin{figure}[H]
    \centering
    \includegraphics[width=0.7\textwidth]{class_distribution.png}
    \caption{CIFAR-10数据集类别分布}
    \label{fig:class_distribution}
\end{figure}

\subsubsection{样本可视化}

图~\ref{fig:sample_images}展示了数据集中各类别的代表性样本。可以观察到，由于图像分辨率较低，某些类别之间的区分度确实较小。

\begin{figure}[H]
    \centering
    \includegraphics[width=0.9\textwidth]{sample_images.png}
    \caption{CIFAR-10各类别样本示例}
    \label{fig:sample_images}
\end{figure}

\subsubsection{像素统计}

我们对训练集的像素值进行了统计分析，计算了每个通道的均值和标准差：
\begin{itemize}[leftmargin=2em]
    \item \textbf{均值:} $[0.4914, 0.4822, 0.4465]$
    \item \textbf{标准差:} $[0.2023, 0.1994, 0.2010]$
\end{itemize}

这些统计量被用于数据标准化，使输入数据具有零均值和单位方差，有助于加速模型收敛。

\subsection{数据划分}

我们将原始的50,000张训练图像进一步划分为训练集和验证集：
\begin{itemize}[leftmargin=2em]
    \item \textbf{训练集:} 45,000张图像 (90\%)
    \item \textbf{验证集:} 5,000张图像 (10\%)
    \item \textbf{测试集:} 10,000张图像 (保持不变)
\end{itemize}

划分采用\textbf{分层抽样}(Stratified Sampling)策略，确保每个子集中各类别的比例与原始数据集一致。这样可以保证验证集能够准确反映模型在各类别上的泛化性能。

\subsection{数据预处理与增强}

\subsubsection{基础预处理}

所有图像都经过以下基础预处理：
\begin{enumerate}[leftmargin=2em]
    \item \textbf{ToTensor:} 将PIL图像转换为PyTorch张量，并将像素值从$[0, 255]$缩放到$[0.0, 1.0]$
    \item \textbf{Normalize:} 使用计算得到的均值和标准差进行标准化
    \begin{equation}
        x_{\text{norm}} = \frac{x - \mu}{\sigma}
    \end{equation}
\end{enumerate}

\subsubsection{数据增强}

为了提高模型的泛化能力，我们在训练阶段应用了以下数据增强技术：

\begin{lstlisting}[caption={数据增强配置},label={lst:augmentation}]
transform_train = transforms.Compose([
    transforms.RandomCrop(32, padding=4),
    transforms.RandomHorizontalFlip(p=0.5),
    transforms.ToTensor(),
    transforms.Normalize(
        mean=[0.4914, 0.4822, 0.4465],
        std=[0.2023, 0.1994, 0.2010]
    ),
])
\end{lstlisting}

\begin{itemize}[leftmargin=2em]
    \item \textbf{RandomCrop:} 在图像四周填充4个像素的0值，然后随机裁剪出$32 \times 32$的区域。这相当于对图像进行了随机平移，增加了位置不变性。
    \item \textbf{RandomHorizontalFlip:} 以0.5的概率对图像进行水平翻转，增加对左右方向的不变性。
    \item \textbf{Normalize:} 使用数据集统计量进行标准化。
\end{itemize}

测试时仅应用ToTensor和Normalize，不使用随机增强，以保证评估的一致性。

% ==================== 第3章 ====================
\section{模型构建}

本节详细介绍三种不同的神经网络架构：多层感知机(MLP)、卷积神经网络(CNN)和ResNet-18。

\subsection{多层感知机 (MLP)}

\subsubsection{模型架构}

MLP是最基础的神经网络结构，我们将$32 \times 32 \times 3 = 3072$维的图像向量展平后输入全连接网络。

\begin{lstlisting}[caption={MLP模型定义},label={lst:mlp}]
class MLP(nn.Module):
    def __init__(self, num_classes=10):
        super(MLP, self).__init__()
        self.flatten = nn.Flatten()
        self.fc1 = nn.Linear(32 * 32 * 3, 1024)
        self.bn1 = nn.BatchNorm1d(1024)
        self.dropout1 = nn.Dropout(0.3)
        self.fc2 = nn.Linear(1024, 512)
        self.bn2 = nn.BatchNorm1d(512)
        self.dropout2 = nn.Dropout(0.3)
        self.fc3 = nn.Linear(512, num_classes)
        
    def forward(self, x):
        x = self.flatten(x)
        x = F.relu(self.bn1(self.fc1(x)))
        x = self.dropout1(x)
        x = F.relu(self.bn2(self.fc2(x)))
        x = self.dropout2(x)
        x = self.fc3(x)
        return x
\end{lstlisting}

\subsubsection{模型特点}

\begin{itemize}[leftmargin=2em]
    \item \textbf{3层全连接结构:} 输入层(3072) $\rightarrow$ 隐藏层1(1024) $\rightarrow$ 隐藏层2(512) $\rightarrow$ 输出层(10)
    \item \textbf{Batch Normalization:} 在每个隐藏层后应用，加速训练并提高稳定性
    \item \textbf{Dropout:} 丢弃率0.3，防止过拟合
    \item \textbf{参数量:} 约380万参数
    \item \textbf{局限性:} 忽略了图像的空间结构信息，将图像视为一维向量
\end{itemize}

\subsection{卷积神经网络 (CNN)}

\subsubsection{模型架构}

CNN通过卷积操作自动提取图像的局部特征，保留了空间结构信息。

\begin{lstlisting}[caption={CNN模型定义},label={lst:cnn}]
class CNN(nn.Module):
    def __init__(self, num_classes=10, use_batchnorm=True, dropout=0.3):
        super(CNN, self).__init__()
        self.use_batchnorm = use_batchnorm
        
        # Block 1
        self.conv1 = nn.Conv2d(3, 64, 3, padding=1)
        self.bn1 = nn.BatchNorm2d(64) if use_batchnorm else nn.Identity()
        self.conv2 = nn.Conv2d(64, 64, 3, padding=1)
        self.bn2 = nn.BatchNorm2d(64) if use_batchnorm else nn.Identity()
        self.pool1 = nn.MaxPool2d(2, 2)
        
        # Block 2
        self.conv3 = nn.Conv2d(64, 128, 3, padding=1)
        self.bn3 = nn.BatchNorm2d(128) if use_batchnorm else nn.Identity()
        self.conv4 = nn.Conv2d(128, 128, 3, padding=1)
        self.bn4 = nn.BatchNorm2d(128) if use_batchnorm else nn.Identity()
        self.pool2 = nn.MaxPool2d(2, 2)
        
        # Block 3
        self.conv5 = nn.Conv2d(128, 256, 3, padding=1)
        self.bn5 = nn.BatchNorm2d(256) if use_batchnorm else nn.Identity()
        self.conv6 = nn.Conv2d(256, 256, 3, padding=1)
        self.bn6 = nn.BatchNorm2d(256) if use_batchnorm else nn.Identity()
        self.pool3 = nn.MaxPool2d(2, 2)
        
        # Block 4
        self.conv7 = nn.Conv2d(256, 512, 3, padding=1)
        self.bn7 = nn.BatchNorm2d(512) if use_batchnorm else nn.Identity()
        self.conv8 = nn.Conv2d(512, 512, 3, padding=1)
        self.bn8 = nn.BatchNorm2d(512) if use_batchnorm else nn.Identity()
        
        # Global Average Pooling + Classifier
        self.gap = nn.AdaptiveAvgPool2d(1)
        self.dropout = nn.Dropout(dropout)
        self.fc = nn.Linear(512, num_classes)
        
    def forward(self, x):
        # Block 1
        x = F.relu(self.bn1(self.conv1(x)))
        x = F.relu(self.bn2(self.conv2(x)))
        x = self.pool1(x)
        
        # Block 2
        x = F.relu(self.bn3(self.conv3(x)))
        x = F.relu(self.bn4(self.conv4(x)))
        x = self.pool2(x)
        
        # Block 3
        x = F.relu(self.bn5(self.conv5(x)))
        x = F.relu(self.bn6(self.conv6(x)))
        x = self.pool3(x)
        
        # Block 4
        x = F.relu(self.bn7(self.conv7(x)))
        x = F.relu(self.bn8(self.conv8(x)))
        
        # GAP and classifier
        x = self.gap(x)
        x = x.view(x.size(0), -1)
        x = self.dropout(x)
        x = self.fc(x)
        return x
\end{lstlisting}

\subsubsection{模型特点}

\begin{itemize}[leftmargin=2em]
    \item \textbf{4个卷积块:} 通道数从64逐步增加到512，提取从低级到高级的特征
    \item \textbf{Batch Normalization:} 每个卷积层后都应用，稳定训练过程
    \item \textbf{Max Pooling:} 降低特征图尺寸，增加感受野，减少计算量
    \item \textbf{Global Average Pooling (GAP):} 替代全连接层，大幅减少参数量，降低过拟合风险
    \item \textbf{Dropout:} 在GAP后应用，丢弃率可配置(默认0.3)
    \item \textbf{参数量:} 约469万参数
\end{itemize}

\subsection{ResNet-18}

\subsubsection{残差学习框架}

ResNet(Residual Network)\cite{he2016deep}通过引入\textbf{残差连接}(Skip Connection)解决了深层网络的梯度消失问题。其核心思想是学习残差映射$\mathcal{F}(x) = \mathcal{H}(x) - x$，而非直接学习期望映射$\mathcal{H}(x)$。

\begin{equation}
    y = \mathcal{F}(x, \{W_i\}) + x
\end{equation}

其中$\mathcal{F}$表示残差映射，$x$是输入，$y$是输出。

\subsubsection{Pre-activation结构}

我们采用了Pre-activation ResNet结构，即在卷积之前先进行Batch Normalization和ReLU激活：

\begin{equation}
    y = \mathcal{F}(\text{BN}(\text{ReLU}(x))) + x
\end{equation}

这种结构具有更好的优化特性\cite{he2016identity}。

\subsubsection{模型架构}

\begin{lstlisting}[caption={ResNet-18基本模块},label={lst:resnet}]
class BasicBlock(nn.Module):
    expansion = 1
    
    def __init__(self, in_planes, planes, stride=1):
        super(BasicBlock, self).__init__()
        self.bn1 = nn.BatchNorm2d(in_planes)
        self.conv1 = nn.Conv2d(in_planes, planes, 3, 
                                stride=stride, padding=1, bias=False)
        self.bn2 = nn.BatchNorm2d(planes)
        self.conv2 = nn.Conv2d(planes, planes, 3, 
                                stride=1, padding=1, bias=False)
        
        self.shortcut = nn.Sequential()
        if stride != 1 or in_planes != self.expansion * planes:
            self.shortcut = nn.Sequential(
                nn.Conv2d(in_planes, self.expansion * planes, 1,
                         stride=stride, bias=False),
                nn.BatchNorm2d(self.expansion * planes)
            )
            
    def forward(self, x):
        out = self.conv1(F.relu(self.bn1(x)))
        out = self.conv2(F.relu(self.bn2(out)))
        out += self.shortcut(x)
        return out
\end{lstlisting}

\subsubsection{模型特点}

\begin{itemize}[leftmargin=2em]
    \item \textbf{18层深度:} 包含17个卷积层和1个全连接层
    \item \textbf{残差连接:} 每个基本块都包含跳跃连接，缓解梯度消失
    \item \textbf{Pre-activation:} BN-ReLU-Conv的顺序，优化训练动态
    \item \textbf{参数量:} 约1,117万参数
    \item \textbf{优势:} 可以训练更深的网络，提取更丰富的特征表示
\end{itemize}

\subsection{训练配置}

三种模型采用统一的训练配置，以确保实验结果的可比性：

\begin{table}[H]
    \centering
    \caption{统一训练配置}
    \label{tab:training_config}
    \begin{tabular}{ll}
        \toprule
        \textbf{配置项} & \textbf{设置值} \\
        \midrule
        优化器 & SGD \\
        学习率 & 0.1 \\
        动量 & 0.9 \\
        权重衰减 & $5 \times 10^{-4}$ \\
        学习率调度 & CosineAnnealingLR \\
        Batch Size & 128 \\
        训练轮数 & 200 (ResNet-18) / 50 (MLP, CNN) \\
        损失函数 & CrossEntropyLoss \\
        \bottomrule
    \end{tabular}
\end{table}

% ==================== 第4章 ====================
\section{实验结果与分析}

\subsection{模型对比实验}

我们首先对比了三种模型在CIFAR-10测试集上的性能。

\subsubsection{准确率对比}

\begin{table}[H]
    \centering
    \caption{三种模型性能对比}
    \label{tab:model_comparison}
    \begin{tabular}{lcc}
        \toprule
        \textbf{模型} & \textbf{测试准确率} & \textbf{参数量} \\
        \midrule
        MLP & 39.31\% & $\sim$3.8M \\
        CNN & 83.66\% & $\sim$4.69M \\
        ResNet-18 & 95.16\% & $\sim$11.17M \\
        \bottomrule
    \end{tabular}
\end{table}

\subsubsection{结果分析}

从表~\ref{tab:model_comparison}可以得出以下结论：

\begin{enumerate}[leftmargin=2em]
    \item \textbf{MLP表现最差:} 仅达到39.31\%的准确率。这是因为MLP将图像展平为一维向量，完全丢失了空间结构信息，无法有效提取局部特征。
    
    \item \textbf{CNN显著提升:} CNN通过卷积操作提取局部特征，并通过层次化结构逐步抽象，准确率提升至83.66\%，比MLP提高了44.35个百分点。
    
    \item \textbf{ResNet-18最优:} 引入残差连接后，网络可以训练得更深，提取更丰富的特征，最终达到95.16\%的准确率，比CNN提高了11.5个百分点。
\end{enumerate}

\subsubsection{训练曲线}

图~\ref{fig:model_curves}展示了三种模型的训练过程。

\begin{figure}[H]
    \centering
    \includegraphics[width=0.9\textwidth]{model_comparison_curves.png}
    \caption{三种模型训练曲线对比}
    \label{fig:model_curves}
\end{figure}

\begin{figure}[H]
    \centering
    \begin{subfigure}[b]{0.32\textwidth}
        \includegraphics[width=\textwidth]{mlp_training_curves.png}
        \caption{MLP训练曲线}
    \end{subfigure}
    \hfill
    \begin{subfigure}[b]{0.32\textwidth}
        \includegraphics[width=\textwidth]{cnn_training_curves.png}
        \caption{CNN训练曲线}
    \end{subfigure}
    \hfill
    \begin{subfigure}[b]{0.32\textwidth}
        \includegraphics[width=\textwidth]{resnet18_training_curves.png}
        \caption{ResNet-18训练曲线}
    \end{subfigure}
    \caption{各模型详细训练曲线}
    \label{fig:individual_curves}
\end{figure}

\subsubsection{混淆矩阵分析}

\begin{figure}[H]
    \centering
    \begin{subfigure}[b]{0.32\textwidth}
        \includegraphics[width=\textwidth]{mlp_confusion_matrix.png}
        \caption{MLP混淆矩阵}
    \end{subfigure}
    \hfill
    \begin{subfigure}[b]{0.32\textwidth}
        \includegraphics[width=\textwidth]{cnn_confusion_matrix.png}
        \caption{CNN混淆矩阵}
    \end{subfigure}
    \hfill
    \begin{subfigure}[b]{0.32\textwidth}
        \includegraphics[width=\textwidth]{resnet18_confusion_matrix.png}
        \caption{ResNet-18混淆矩阵}
    \end{subfigure}
    \caption{各模型混淆矩阵对比}
    \label{fig:confusion_matrices}
\end{figure}

从混淆矩阵可以看出，ResNet-18的对角线元素明显更突出，说明其分类置信度更高。MLP则存在大量类间混淆，特别是猫和狗、鹿和马等相似类别之间。

\subsection{消融实验}

为了深入理解各技术组件的作用，我们在CNN模型(训练50轮)上进行了系统的消融实验。

\subsubsection{Batch Normalization的影响}

\begin{table}[H]
    \centering
    \caption{Batch Normalization消融实验}
    \label{tab:ablation_bn}
    \begin{tabular}{lc}
        \toprule
        \textbf{配置} & \textbf{验证准确率} \\
        \midrule
        使用BN & 89.88\% \\
        不使用BN & 10.00\% \\
        \bottomrule
    \end{tabular}
\end{table}

\textbf{分析:} 不使用Batch Normalization时，模型完全无法学习，验证准确率仅10\%(相当于随机猜测)。这是因为：
\begin{itemize}[leftmargin=2em]
    \item 学习率0.1对于没有BN的网络过大，导致梯度爆炸
    \item 内部协变量偏移(Internal Covariate Shift)使深层网络难以训练
    \item BN通过归一化层输入，允许使用更大的学习率，加速收敛
\end{itemize}

\textbf{结论:} Batch Normalization对于使用较大学习率(如0.1)和SGD优化器训练CNN至关重要。

\begin{figure}[H]
    \centering
    \includegraphics[width=0.7\textwidth]{ablation_batchnorm.png}
    \caption{Batch Normalization消融实验结果}
    \label{fig:ablation_bn}
\end{figure}

\subsubsection{Dropout的影响}

\begin{table}[H]
    \centering
    \caption{Dropout消融实验}
    \label{tab:ablation_dropout}
    \begin{tabular}{lc}
        \toprule
        \textbf{Dropout率} & \textbf{验证准确率} \\
        \midrule
        0.0 & 90.62\% \\
        0.3 & 89.58\% \\
        0.5 & 88.08\% \\
        \bottomrule
    \end{tabular}
\end{table}

\textbf{分析:} 有趣的是，不使用Dropout(0.0)反而获得了最高的验证准确率(90.62\%)。这表明：
\begin{itemize}[leftmargin=2em]
    \item 在仅训练50轮的情况下，模型尚未充分过拟合
    \item Dropout作为一种正则化技术，在模型未过拟合时可能阻碍学习
    \item 随着训练轮数增加，Dropout的作用可能会显现
\end{itemize}

\textbf{结论:} 在50轮训练设置下，模型欠拟合而非过拟合，Dropout未能带来性能提升。

\begin{figure}[H]
    \centering
    \includegraphics[width=0.7\textwidth]{ablation_dropout.png}
    \caption{Dropout消融实验结果}
    \label{fig:ablation_dropout}
\end{figure}

\subsubsection{优化器选择}

\begin{table}[H]
    \centering
    \caption{优化器对比实验}
    \label{tab:ablation_optimizer}
    \begin{tabular}{lc}
        \toprule
        \textbf{优化器} & \textbf{验证准确率} \\
        \midrule
        SGD (lr=0.1) & 89.58\% \\
        Adam (lr=0.001) & 89.52\% \\
        \bottomrule
    \end{tabular}
\end{table}

\textbf{分析:} 两种优化器在适当的学习率设置下表现相当：
\begin{itemize}[leftmargin=2em]
    \item SGD配合较大的学习率(0.1)和动量(0.9)表现良好
    \item Adam使用自适应学习率，对学习率初始值不那么敏感
    \item 两者差距仅0.06\%，在实验误差范围内
\end{itemize}

\textbf{结论:} 优化器选择对最终性能影响较小，关键在于选择合适的学习率。

\begin{figure}[H]
    \centering
    \includegraphics[width=0.7\textwidth]{ablation_optimizer.png}
    \caption{优化器对比实验结果}
    \label{fig:ablation_optimizer}
\end{figure}

\subsubsection{学习率调度策略}

\begin{table}[H]
    \centering
    \caption{学习率调度策略对比}
    \label{tab:ablation_scheduler}
    \begin{tabular}{lc}
        \toprule
        \textbf{调度策略} & \textbf{验证准确率} \\
        \midrule
        CosineAnnealing & 89.80\% \\
        ReduceLROnPlateau & 86.76\% \\
        无调度 & 80.54\% \\
        \bottomrule
    \end{tabular}
\end{table}

\textbf{分析:} 学习率调度对性能有显著影响：
\begin{itemize}[leftmargin=2em]
    \item \textbf{CosineAnnealing:} 按照余弦函数平滑降低学习率，使模型能够充分收敛到损失平面的低谷
    \item \textbf{ReduceLROnPlateau:} 在验证损失停滞时降低学习率，但可能过于保守
    \item \textbf{无调度:} 固定学习率导致后期无法精细调整，性能最差
\end{itemize}

\textbf{结论:} CosineAnnealing调度策略表现最佳，建议作为默认选择。

\begin{figure}[H]
    \centering
    \includegraphics[width=0.7\textwidth]{ablation_scheduler.png}
    \caption{学习率调度策略对比}
    \label{fig:ablation_scheduler}
\end{figure}

\subsubsection{数据增强策略}

\begin{table}[H]
    \centering
    \caption{数据增强消融实验}
    \label{tab:ablation_aug}
    \begin{tabular}{lc}
        \toprule
        \textbf{增强策略} & \textbf{验证准确率} \\
        \midrule
        完整增强(RandomCrop+HFlip+Normalize) & 89.38\% \\
        仅Normalize & 84.04\% \\
        无增强 & 79.38\% \\
        \bottomrule
    \end{tabular}
\end{table}

\textbf{分析:} 数据增强对泛化能力有重要影响：
\begin{itemize}[leftmargin=2em]
    \item 完整增强策略比无增强提高了10个百分点
    \item RandomCrop和HorizontalFlip通过增加训练样本的多样性，有效提升了模型的泛化能力
    \item 即使仅使用Normalize也能获得一定提升，因为标准化使优化 landscape 更平滑
\end{itemize}

\textbf{结论:} 数据增强是提升模型泛化能力的有效手段，建议始终使用。

\begin{figure}[H]
    \centering
    \includegraphics[width=0.7\textwidth]{ablation_data_aug.png}
    \caption{数据增强消融实验结果}
    \label{fig:ablation_aug}
\end{figure}

\subsubsection{消融实验总结}

\begin{figure}[H]
    \centering
    \includegraphics[width=0.9\textwidth]{ablation_comparison.png}
    \caption{各消融实验验证准确率对比}
    \label{fig:ablation_summary}
\end{figure}

综合所有消融实验，各技术组件的重要性排序为：
\begin{equation}
    \text{BatchNorm} > \text{Scheduler} > \text{Data Aug} > \text{Dropout} \approx \text{Optimizer}
\end{equation}

\subsection{ResNet-18增强训练}

为了进一步提升ResNet-18的性能，我们尝试了Label Smoothing和Warmup技术。

\subsubsection{Label Smoothing}

Label Smoothing\cite{szegedy2016rethinking}通过将one-hot标签替换为软标签来防止模型过于自信：

\begin{equation}
    q(k|x) = (1 - \epsilon) \delta_{k,y} + \epsilon / K
\end{equation}

其中$\epsilon = 0.1$是平滑参数，$K = 10$是类别数。

\subsubsection{Warmup策略}

训练初期使用线性Warmup策略，在前5个epoch内将学习率从0逐渐增加到初始值0.1：

\begin{equation}
    lr_t = lr_{\text{initial}} \times \frac{t}{T_{\text{warmup}}}, \quad t \in [1, T_{\text{warmup}}]
\end{equation}

\subsubsection{实验结果}

\begin{table}[H]
    \centering
    \caption{ResNet-18增强训练结果}
    \label{tab:resnet_enhanced}
    \begin{tabular}{lc}
        \toprule
        \textbf{配置} & \textbf{验证准确率} \\
        \midrule
        基线(ResNet-18) & 95.16\% \\
        + Label Smoothing (0.1) + Warmup (5 epochs) & 95.16\% \\
        \bottomrule
    \end{tabular}
\end{table}

\textbf{分析:} 在此实验设置下，Label Smoothing和Warmup并未带来额外提升：
\begin{itemize}[leftmargin=2em]
    \item 基线模型已经达到较高性能，提升空间有限
    \item 200轮训练配合CosineAnnealing已经使模型充分收敛
    \item 这些技术可能在更困难的分类任务或更长的训练周期中显现价值
\end{itemize}

\textbf{结论:} 对于CIFAR-10和ResNet-18的组合，基线配置已经足够优秀。

\begin{figure}[H]
    \centering
    \includegraphics[width=0.8\textwidth]{resnet18_smooth.png}
    \caption{ResNet-18增强训练曲线}
    \label{fig:resnet_enhanced}
\end{figure}

\subsection{各类别性能分析}

\begin{figure}[H]
    \centering
    \includegraphics[width=0.9\textwidth]{per_class_comparison.png}
    \caption{各模型每类准确率对比}
    \label{fig:per_class}
\end{figure}

从图~\ref{fig:per_class}可以看出：
\begin{itemize}[leftmargin=2em]
    \item ResNet-18在所有类别上都显著优于其他模型
    \item 飞机、汽车、船等具有明显轮廓的物体更容易分类
    \item 猫、狗、鹿等动物类别更具挑战性，类间相似性高
    \item MLP在动物类别上表现尤其差，说明其无法提取区分性特征
\end{itemize}

% ==================== 第5章 ====================
\section{结论与可能改进}

\subsection{主要发现}

通过本次大作业的系统性实验，我们得出以下关键结论：

\begin{enumerate}[leftmargin=2em]
    \item \textbf{网络架构至关重要:} ResNet-18(95.16\%)显著优于CNN(83.66\%)和MLP(39.31\%)。残差连接使深层网络训练成为可能，能够学习更复杂的特征表示。
    
    \item \textbf{Batch Normalization不可或缺:} 消融实验表明，没有BN的网络完全无法训练(10\% vs 89.88\%)。BN不仅加速收敛，还使大学习率训练成为可能。
    
    \item \textbf{学习率调度策略重要:} CosineAnnealing比无调度提高了9.26个百分点，是仅次于BN的重要技术。
    
    \item \textbf{数据增强提升泛化:} 完整的数据增强策略比无增强提高了10个百分点，是防止过拟合的有效手段。
    
    \item \textbf{Dropout作用有限:} 在50轮训练设置下，模型尚未过拟合，Dropout未能带来提升。这可能需要更长的训练周期才能显现价值。
    
    \item \textbf{优化器选择相对次要:} SGD和Adam在适当学习率下表现相当，差异在实验误差范围内。
\end{enumerate}

\subsection{可能的改进方向}

基于本次实验的经验，我们提出以下可能的改进方向：

\subsubsection{1. 更强的数据增强}

\begin{itemize}[leftmargin=2em]
    \item \textbf{Cutout:} 随机遮挡图像的一部分，强迫模型学习更鲁棒的特征
    \item \textbf{Mixup:} 将两张图像按一定比例混合，生成新的训练样本，提高泛化能力
    \item \textbf{AutoAugment:} 使用强化学习自动搜索最优增强策略
    \item \textbf{RandAugment:} 简化的自动增强方法，无需额外搜索成本
\end{itemize}

\subsubsection{2. 更深的网络架构}

\begin{itemize}[leftmargin=2em]
    \item \textbf{ResNet-34/50:} 尝试更深的残差网络，可能提取更丰富的特征
    \item \textbf{Wide ResNet:} 增加通道数而非深度，在CIFAR-10上表现优异
    \item \textbf{DenseNet:} 密集连接，特征重用，参数效率更高
    \item \textbf{EfficientNet:} 复合缩放策略，平衡深度、宽度和分辨率
\end{itemize}

\subsubsection{3. 超参数优化}

\begin{itemize}[leftmargin=2em]
    \item \textbf{学习率搜索:} 使用网格搜索或随机搜索寻找最优初始学习率
    \item \textbf{权重衰减调优:} 尝试不同的权重衰减系数($10^{-3}$到$10^{-5}$)
    \item \textbf{Batch Size影响:} 探索不同Batch Size对性能和收敛速度的影响
    \item \textbf{贝叶斯优化:} 使用更高效的超参数优化算法
\end{itemize}

\subsubsection{4. 模型集成}

\begin{itemize}[leftmargin=2em]
    \item \textbf{多模型集成:} 训练多个独立模型，取预测平均
    \item \textbf{快照集成(Snapshot Ensembles):} 利用单个模型的多个检查点
    \item \textbf{知识蒸馏:} 使用大模型(ResNet-50)指导小模型训练
\end{itemize}

\subsubsection{5. 延长训练周期}

\begin{itemize}[leftmargin=2em]
    \item \textbf{更长的训练:} 尝试400-600轮训练，配合强正则化
    \item \textbf{渐进式调整:} 逐步增加图像分辨率或网络深度
    \item \textbf{测试时增强(TTA):} 对测试图像应用多种增强，取平均预测
\end{itemize}

\subsection{总结}

本报告通过系统的实验设计和详细的消融分析，深入探讨了深度学习图像分类中的关键技术。从简单的MLP到先进的ResNet-18，我们见证了网络架构演进带来的巨大性能提升。消融实验揭示了Batch Normalization、学习率调度和数据增强等技术的实际价值，为今后的深度学习实践提供了宝贵的经验。

CIFAR-10作为一个经典的基准数据集，虽然图像分辨率较低，但涵盖了丰富的视觉概念和类间挑战。在这个数据集上获得的洞察，可以迁移到更复杂、更实用的计算机视觉任务中。

% ==================== 参考文献 ====================
\newpage
\printbibliography

\end{document}
